\begin{abstract}
La distribution des galaxies contient plus d'information cosmologique que ne peuvent en extraire les statistiques usuelles à deux points : la toile cosmique non linéaire de nœuds, de filaments, de feuillets et de vides est elle aussi sensible à la croissance des structures, au contenu en matière, à l'énergie noire et aux neutrinos. Nous cherchons à savoir quelle part de cette information supplémentaire peut être récupérée avec ASTRA, un algorithme qui attribue un environnement à chaque point à partir de ses voisins dans une tessellation de Delaunay des galaxies et d'un catalogue aléatoire, sans aucune échelle de lissage. Nous séparons les galaxies et les points aléatoires selon leur environnement ASTRA, mesurons leurs auto- et cross-corrélations, et les comparons à la fonction à deux points des galaxies et au clustering par découpage en densité (density split), la statistique d'environnement de référence. Nous construisons les covariances et les dérivées cosmologiques nécessaires aux prévisions de Fisher à partir de trois suites de simulations : des mocks de galaxies de type DESI couvrant 85 cosmologies (AbacusSummit), 1786 mocks indépendants du programme officiel de DESI, et les simulations Quijote. En chemin, nous entraînons un émulateur par réseau de neurones de toutes les statistiques et l'utilisons pour un test d'inférence de bout en bout. Les statistiques d'environnement ASTRA contraignent $\omega_b$, $\omega_{\rm cdm}$, $n_s$, $\sigma_8$, $w_0$ et $w_a$ environ deux fois mieux que la fonction à deux points. Elles portent autant d'information que le density split, et les deux sont complémentaires : les combiner améliore encore les contraintes. Les classes de toile cosmique originales d'ASTRA, testées sur Quijote, contraignent $\sigma_8$ deux à trois fois mieux que le density split. Nous décrivons aussi les principales difficultés de ce type de prévision (estimer des dérivées à partir de suites d'émulation, choisir une covariance, traiter des dérivées bruitées) et la manière dont nous les avons traitées.
\end{abstract}
